% !TEX root = task-1-14-5.tex
% Editable task bank. Append new BankTask blocks at the end of the document.
\documentclass[12pt,a4paper]{article}
\usepackage{physics-textbook}
\usepackage{graphicx,needspace}
\usetikzlibrary{decorations.pathmorphing}
\newif\ifShowSolutions
\ifdefined\TutorEdition
  \ShowSolutionstrue
\else
  \ShowSolutionsfalse % Change to \ShowSolutionstrue to display solutions.
\fi
\newcounter{banktask}
\newsavebox{\BankTaskBox}
\newcommand{\BankTaskContent}[3]{%
  \noindent\textbf{\thebanktask.}\quad #2\par
  \ifShowSolutions
    \Needspace{50mm}% Keep the solution heading with its opening text or diagram.
    #3\par
    \ifdefined\TutorEdition
      \ifcsname TutorSource#1\endcsname
        \nopagebreak\smallskip{\small\csname TutorSource#1\endcsname\par}%
      \fi
    \fi
  \fi
}
\newcommand{\BankTask}[3]{%
  \refstepcounter{banktask}\par\medskip
  \sbox{\BankTaskBox}{\begin{minipage}{\linewidth}%
    \BankTaskContent{#1}{#2}{#3}\end{minipage}}%
  \ifdim\dimexpr\ht\BankTaskBox+\dp\BankTaskBox\relax>\dimexpr\textheight-20mm\relax
    \Needspace{35mm}\BankTaskContent{#1}{#2}{#3}%
  \else
    \Needspace{\dimexpr\ht\BankTaskBox+\dp\BankTaskBox\relax}%
    \noindent\usebox{\BankTaskBox}\par
  \fi
  \medskip
}
\begin{document}
\BookSubsection{1.14.5}{Закон сохранения проекции импульса}


% Source sdamgia-3351: https://phys-ege.sdamgia.ru/problem?id=3351
\BankTask{sdamgia-3351}{
Камень массой $m_1=4\,\text{кг}$ падает под углом $60^\circ$ к горизонту со скоростью $10\,\text{м/с}$ в тележку с песком, покоящуюся на горизонтальных рельсах (см. рис.). Чему равен импульс тележки с песком и камнем после падения камня?


\begin{center}\begin{tikzpicture}[task graph]
\draw[task curve] (0,0)--(3,0);\draw[fill=orange!25] (.7,.5) rectangle (2.5,1.05);\draw[task curve] (.7,1.05)..controls(1.1,1.3)and(1.7,.9)..(2.5,1.05);
\draw (.95,.3) circle (.2);\draw (2.25,.3) circle (.2);\node at (1.6,.8) {$m_2$};
\draw[fill=orange!25] (.5,2.9) circle (.18) node[left] {$m_1$};\draw[task curve] (.7,2.9)--(2.5,2.9);\draw[task curve,->] (.6,2.72)--(1.5,1.2);\draw (1.2,2.9) arc (0:-60:.7) node[midway,right] {$\alpha$};
\end{tikzpicture}\end{center}
}{
\smallskip\noindent\textbf{Решение}\par
Поскольку на систему в горизонтальном направлении не действует никаких внешних сил, в этом направлении выполняется закон сохранения импульса. Полный горизонтальный импульс системы до столкновения должен быть равен полному горизонтальному импульсу системы после столкновения. Таким образом, искомый импульс тележки с песком и камнем после падения равен:

$p=m_1v_{\text{камн}}\cos\alpha=4\,\text{кг}\cdot10\,\text{м/с}\cdot\cos60^\circ=20\,\text{кг}\cdot\text{м/с}.$

Ответ: $20$.
}

% Source sdamgia-24944: https://phys-ege.sdamgia.ru/problem?id=24944
\BankTask{sdamgia-24944}{
Из незакрепленной пушки стреляют снарядом массой $20\,\text{кг}$, который вылетает из ствола в горизонтальном направлении со скоростью $102\,\text{м/с}$ относительно пушки. Пушка при этом откатывается, приобретая относительно земли скорость $2\,\text{м/с}$. Чему равна масса пушки, если массой сгоревшего порохового заряда можно пренебречь?
}{
\smallskip\noindent\textbf{Решение}\par
По закону сохранения импульса $0=m_1\vec v_1+m_2\vec v_2.$ Причем, скорость снаряда относительно земли по правилу сложения скоростей равна $\vec v_1=\vec v_{\text{отн}}+\vec v_2.$ Тогда получаем: $0=m_1(\vec v_{\text{отн}}+\vec v_2)+m_2\vec v_2.$ Учитывая, что пушка откатится в сторону, противоположную направлению движения снаряда, получаем проекцию данного уравнения на горизонтальную ось: $0=m_1v_{\text{отн}}-m_1v_2-m_2v_2,$ откуда получаем искомую величину: 

$m_2=\frac{m_1(v_{\text{отн}}-v_2)}{v_2}=\frac{20\cdot(102-2)}2=1000\,\text{кг}.$

Ответ: $1000$.
}

% Source sdamgia-8934: https://phys-ege.sdamgia.ru/problem?id=8934
\BankTask{sdamgia-8934}{
Тележка движется по инерции по гладким горизонтальным рельсам со скоростью $4\,\text{м/с}$. На тележку вертикально сверху аккуратно опускают мешочек с песком. Масса мешочка в $3$ раза меньше массы тележки. Чему будет равен модуль скорости тележки с мешочком после того, как проскальзывание мешочка относительно тележки прекратится?
}{
\smallskip\noindent\textbf{Решение}\par
Запишем закон сохранения импульса:

$3mV_1=(3m+m)V_2,$

где $m$ — масса мешочка с песком. Найдем отсюда скорость тележки с мешочком:

$V_2=\frac34V_1=3\,\text{м/с}.$

Ответ: $3$.
}

% Source sdamgia-42029: https://phys-ege.sdamgia.ru/problem?id=42029
\BankTask{sdamgia-42029}{
Камень массой $m=2\,\text{кг}$ падает под углом $\alpha=30^\circ$ к вертикали со скоростью $v=15\,\text{м/с}$ в ящик с песком общей массой $M=18\,\text{кг}$, который покоится на гладком горизонтальном столе (см. рисунок). Определите скорость ящика после застревания камня в песке. 


\begin{center}\begin{tikzpicture}[task graph]
\draw[fill=orange!25] (0,0) rectangle (3,.7);\draw[task curve] (0,.7)..controls(.6,1.1)and(1.8,.6)..(3,.9);\draw[task curve] (0,.7)--(0,1.1) (3,.7)--(3,1.1);
\draw[task curve] (-.4,0)--(3.4,0);\foreach \x in {-.4,-.2,...,3.4}{\draw (\x,0)--++(-.12,-.12);}
\fill (.2,2.7) circle (.19) node[above left] {$m$};\draw[task curve] (.2,2.7)--(.2,1.5);\draw[task curve,->] (.35,2.48)--(1.3,1.3) node[right] {$\vec v$};\draw (.2,2) arc (-90:-52:.7) node[midway,below] {$\alpha$};
\end{tikzpicture}\end{center}
}{
\smallskip\noindent\textbf{Решение}\par
По закону сохранения импульса тел при неупругом взаимодействии

$mv\sin\alpha=(m+M)u,$

откуда скорость ящика после взаимодействия равна

$u=\frac{mv\sin\alpha}{m+M}=\frac{2\cdot15\cdot0{,}5}{2+18}=0{,}75\,\text{м/с}.$

Ответ: $0{,}75$.
}

% Source sdamgia-26714: https://phys-ege.sdamgia.ru/problem?id=26714
\BankTask{sdamgia-26714}{
По гладкой горизонтальной плоскости движутся вдоль осей $x$ и $y$ две шайбы с импульсами по модулю $p_{1}=2\,\text{кг}\cdot\text{м/с}$ и $p_{2}=3{,}5\,\text{кг}\cdot\text{м/с}$ (см. рис.). После их соударения вторая шайба продолжает двигаться по оси $y$ в прежнем направлении. Модуль импульса первой шайбы сразу после удара $p^{\prime}_{1}=2{,}5\,\text{кг}\cdot\text{м/с}$. Найдите модуль импульса второй шайбы сразу после удара. 


\begin{center}\begin{tikzpicture}[task graph]
\draw[task axis] (-1.8,0)--(2,0) node[below] {$x$};\draw[task axis] (0,-1.8)--(0,1.6) node[left] {$y$};
\draw[fill=orange!25] (-1.7,0) circle (.2);\draw[task curve,->] (-1.7,0)--(-.6,0) node[midway,below] {$\vec p_1$};
\draw[fill=orange!25] (0,-1.6) circle (.2);\draw[task curve,->] (0,-1.6)--(0,-.6) node[midway,left] {$\vec p_2$};
\end{tikzpicture}\end{center}
}{
\smallskip\noindent\textbf{Решение}\par
По закону сохранения импульса $\vec p_1+\vec p_2=\vec p_1\,\!{}^{\prime}+\vec p_2\,\!{}^{\prime}.$
В проекции на оси х и у $p_1=p_{1x}^{\prime};$ $p_2=p_{1y}^{\prime}+p_2^{\prime}.$
Тогда $p_2^{\prime}=p_2-p_{1y}^{\prime}.$ Проекцию импульса первого тела на ось у после соударения найдем из теоремы Пифагора:

$p_{1y}^{\prime}=\sqrt{(p_1^{\prime})^2-p_{1x}^2}=\sqrt{2{,}5^2-2^2}=1{,}5$ $\text{кг}\cdot\text{м/с}.$

Отсюда импульс второго тела после взаимодействия равен:

$p_2^{\prime}=p_2-p_{1y}^{\prime}=3{,}5-1{,}5=2$ $\text{кг}\cdot\text{м/с}.$

Ответ: $2$.
}

% Source sdamgia-10633: https://phys-ege.sdamgia.ru/problem?id=10633
\BankTask{sdamgia-10633}{
Навстречу тележке массой $4{,}75\,\text{кг}$, движущейся по инерции равномерно со скоростью $2\,\text{м/с}$ по гладким горизонтальным рельсам, летит шар массой $0{,}25\,\text{кг}$ со скоростью $40\,\text{м/с}$. После столкновения шар застревает в песке, насыпанном на тележку. Определите, во сколько раз отличаются модули начального (до застревания в песке) и конечного импульса шара в системе отсчета, связанной с рельсами.
}{
\smallskip\noindent\textbf{Решение}\par
Запишем закон сохранения импульса в проекции на ось, вдоль которой движутся тележка и шар:

$mv_1-Mv_2=(m+M)V,$

где $m$ и $M$ — массы шара и тележки соответственно, $v_1$ и $v_2$ — скорости шара и тележки до столкновения, $V$ — скорость тележки с шаром после столкновения.

Выразим скорость после столкновения:

$V=\frac{mv_1-Mv_2}{m+M}.$

Тогда отношение начального импульса шара к конечному равно:

$\frac{mv_1}{mV}=\frac{v_1}V=\frac{v_1(m+M)}{mv_1-Mv_2}=\frac{40\cdot(0{,}25+4{,}75)}{0{,}25\cdot40-4{,}75\cdot2}=400.$

Ответ: $400$.
}

% Source sdamgia-741: https://phys-ege.sdamgia.ru/problem?id=741
\BankTask{sdamgia-741}{
По гладкой горизонтальной плоскости по осям $x$ и $y$ движутся две шайбы с импульсами, равными по модулю $p_1=2\,\text{кг}\cdot\text{м/с}$ и $p_2=3{,}5\,\text{кг}\cdot\text{м/с},$ как показано на рисунке. После соударения вторая шайба продолжает двигаться по оси $y$ в прежнем направлении с импульсом, равным по модулю $p_3=2\,\text{кг}\cdot\text{м/с}.$ Найдите модуль импульса первой шайбы после удара.


\begin{center}\begin{tikzpicture}[task graph]
\draw[task axis] (-1.8,0)--(2,0) node[below] {$x$};\draw[task axis] (0,-1.8)--(0,1.6) node[left] {$y$};
\draw[fill=orange!25] (-1.7,0) circle (.2);\draw[task curve,->] (-1.7,0)--(-.6,0) node[midway,below] {$\vec p_1$};
\draw[fill=orange!25] (0,-1.6) circle (.2);\draw[task curve,->] (0,-1.6)--(0,-.6) node[midway,left] {$\vec p_2$};
\end{tikzpicture}\end{center}
}{
\smallskip\noindent\textbf{Решение}\par
Поскольку на шайбы не действует никаких внешних сил, для них выполняется закон сохранения импульса:

$\vec p_1+\vec p_2=\vec p_3+\vec p_4.$

Тогда $\vec p_4=\vec p_1+\vec p_2-\vec p_3.$
В проекции на выбранные оси:

$p_{4x}=p_{1x}$
$p_{4y}=p_2-p_3.$

Тогда модуль вектора импульса первой шайбы после соударения равен

$p_4=\sqrt{p_{4x}^2+p_{4y}^2}=\sqrt{p_1^2+(p_2-p_3)^2}=$
$=\sqrt{(2\,\text{кг}\cdot\text{м/с})^2+(3{,}5\,\text{кг}\cdot\text{м/с}-2\,\text{кг}\cdot\text{м/с})^2}=2{,}5\,\text{кг}\cdot\text{м/с}.$

Ответ: $2{,}5\,\text{кг}\cdot\text{м/с}$.
}

% Source sdamgia-38596: https://phys-ege.sdamgia.ru/problem?id=38596
\BankTask{sdamgia-38596}{
Пластилиновый шарик в момент $t=0$ бросают с горизонтальной поверхности земли с некоторой начальной скоростью под углом $\alpha$ = $30^\circ$ к горизонту. Одновременно с некоторой высоты над поверхностью земли начинает падать из состояния покоя другой такой же шарик. Спустя время $\tau$ = $1\,\text{с}$ шарики абсолютно неупруго сталкиваются в воздухе. Сразу после столкновения скорость шариков направлена горизонтально. Какова начальная скорость $v_{0}$ шарика, брошенного под углом к горизонту? Сопротивлением воздуха пренебречь. Какие законы Вы использовали для описания движения? Обоснуйте их применение к данному случаю.
}{
\smallskip\noindent\textbf{Решение}\par
Обоснование
Систему отсчета, связанную с Землей, будем считать инерциальной. Тела движутся поступательно, поэтому их можно принять за материальные точки. По условию задачи силой сопротивления воздуха можно пренебречь. Поэтому тела движутся в поле Земного тяготения. Следовательно, для описания движения тел до столкновения можно использовать законы равноускоренного движения.
Перед столкновением на шарики действуют силы тяжести, которые являются внешними. Поскольку при столкновении на шарики действуют внутренние силы взаимодействия, которые довольно велики из-за кратковременности удара, то по сравнению с ними силой тяжести можно пренебречь, а значит можно пренебречь внешними силами, действующими на тела в процессе соударения. Поэтому можно воспользоваться законом сохранения импульса.

Перейдем к решению. Введем координатные оси: горизонтальную ось $OX$ и направленную вверх ось $OY$. Начало координат разместим в точке старта первого шарика.
До столкновения вертикальные координата и проекции скоростей шариков меняются по законам:

$y_1=v_0\sin\alpha\cdot t-\frac{gt^2}2,$
$v_{y1}=v_0\sin\alpha-gt,$
$v_{y2}=-gt.$
После столкновения импульс слипшихся шариков равен сумме импульсов до столкновения: $\vec p_1+\vec p_2=\vec p_0,$ где $p_1=mv_1,$ $p_2=mv_2.$ 
Исходя из того, что скорость слипшихся шариков горизонтальна, а их массы одинаковы, в проекции на вертикальную ось получаем: $mv_{y1}+mv_{y2}=0,$ откуда $v_{y1}=-v_{y2},$ значит, $v_0\sin\alpha-gt=gt.$
Тогда начальная скорость первого шарика равна

$v_0=\frac{2gt}{\sin\alpha}=\frac{2\cdot10\cdot1}{0{,}5}=40\,\text{м/с}.$

Ответ: $40\,\text{м/с}$.
}

% Source sdamgia-25911: https://phys-ege.sdamgia.ru/problem?id=25911
\BankTask{sdamgia-25911}{
Граната, летевшая с некоторой скоростью, разрывается на две части. Первый осколок летит под углом $90^{\circ}$ к первоначальному направлению со скоростью $40\,\text{м/с}$, а второй  — под углом $30^{\circ}$ со скоростью $20\,\text{м/с}$. Чему равно отношение массы второго осколка к массе первого осколка.
Какие законы Вы используете для описания взрыва гранаты? Обоснуйте их применение к данному случаю.
}{
\smallskip\noindent\textbf{Решение}\par
Обоснование 
Задачу будем решать в инерциальной системе отсчета, связанной с Землей. Все тела будем считать материальными точками. Трением гранаты и осколков о воздух пренебрежем. Для описания разрыва гранаты использован закон сохранения импульса системы тел. Он выполняется в инерциальной системе отсчета, если сумма внешних сил, приложенных к телам системы, равна нулю. В данном случае из-за отсутствия сопротивления воздуха внешней силой является только сила тяжести $mg$, которая не равна нулю. Но этим можно пренебречь, считая время разрыва гранаты малым. За малое время разрыва импульс каждого из осколков меняется на конечную величину за счет больших внутренних сил, разрывающих снаряд при взрыве. По сравнению с этими большими силами конечная сила тяжести пренебрежимо мала. Поскольку время разрыва гранаты считаем малым, то можно пренебречь и изменением потенциальной энергии гранаты и его осколков в поле тяжести в процессе разрыва. В инерциальной системе отсчета выполняется закон сохранения импульса тел.

Перейдем к решению. \begin{center}
\begin{tikzpicture}[task graph,>=stealth]
\coordinate (O) at (0,0);
\coordinate (A) at (0,1.5);
\coordinate (B) at ({1.5*sqrt(3)},0);
\coordinate (C) at ({1.5*sqrt(3)},-1.5);
\draw[dashed] (A)--(B)--(C);
\draw[->] (O)--(A) node[midway,left] {$\vec p_1$};
\draw[->] (O)--(B) node[midway,above] {$\vec p_0$};
\draw[->] (O)--(C) node[midway,below] {$\vec p_2$};
\end{tikzpicture}
\end{center} Поскольку действием внешних сил за время взрыва можно пренебречь, для системы должен выполняться закон сохранения импульса. Импульс ядра до взрыва должен быть равен сумме импульсов осколков после взрыва. Спроектируем закон сохранения импульса на ось, перпендикулярную направлению движения ядра: $0=p_1 - p_2 \sin 30 ^{ \circ } \Longleftrightarrow m_1 v _1=m_2 v _2 \sin 30 ^{ \circ } .$ Следовательно, отношение массы второго осколка к массе первого осколка равно:

\[\frac{m_2}{m_1} = \frac{v _1}{v _2 \sin 30 ^{ \circ }} = \frac{40\,\text{м/с}}{20\,\text{м/с} \cdot \frac{1}{2}} =4.\]

Ответ: $4$.
}


% Source 24378: https://phys-ege.sdamgia.ru/problem?id=24378
\BankTask{24378}{
Сталкиваются и слипаются два разных по массе пластилиновых шарика, причем векторы их скоростей непосредственно перед столкновением направлены навстречу друг другу и одинаковы по модулю: $v_1$ = $v_2$  =  $1\,\text{м/с}$. Во сколько раз масса тяжелого шарика больше, чем легкого, если сразу после столкновения их скорость стала равной (по модулю) $0{,}5\,\text{м/с}$?
}{
\smallskip\noindent\textbf{Решение}\par
По закону сохранения импульса для шариков массами m и M при их неупругом взаимодействии $m\vec v _1 + M\vec v _2= ( M + m ) \vec u.$ Учитывая, что скорости шаров до взаимодействия были равными и направлены противоположно, получаем $M v - m v = ( M + m ) u.$ Отсюда $M(v-u)=m(v+u)$. 

\[\frac{M}{m} = \frac{u + v}{v - u} = \frac{0{,}5 + 1}{1 - 0{,}5} =3.\]

Ответ: в $3$ раза.
}

% Source 39272: https://phys-ege.sdamgia.ru/problem?id=39272
\BankTask{39272}{
Пластилиновый шарик в момент t  =  $0$ бросают с горизонтальной поверхности Земли под углом $\alpha$ к горизонту. Одновременно с некоторой высоты над поверхностью Земли начинает падать из состояния покоя другой такой же шарик. Шарики абсолютно неупруго сталкиваются в воздухе. Сразу после столкновения скорость шариков направлена горизонтально. Время от столкновения шариков до их падения на Землю равно $\tau$. С какой начальной скоростью $\upsilon_{0}$ был брошен первый шарик? Сопротивлением воздуха пренебречь.
}{
\smallskip\noindent\textbf{Решение}\par
\begin{center}\begin{tikzpicture}[task graph,>=stealth]\draw[->](0,0)--(8.8,0)node[right]{$x$};\draw[->](0,0)--(0,4.8)node[above]{$y$};\node[below left]at(0,0){$0$};\draw[domain=0:8.4,samples=60]plot(\x,{1.333333*\x-.15873016*\x*\x});\draw[->](0,0)--(.8,1.066667)node[above left]{$\vec v_0$};\draw(.6,0)arc(0:53.1301:.6)node[midway,right]{$\alpha$};\coordinate(B)at(2.1,2.1);\coordinate(A)at(2.1,3.4);\fill(B)circle(1.7pt)node[left]{$B$};\fill(A)circle(1.7pt)node[left]{$A$};\draw(A)--(B);\draw[->,thick](B)--(3.8,3.233333)node[above]{$m\vec v_1$};\draw[->,thick](B)--(3.8,2.1)node[right]{$2m\vec u_0$};\draw[->,thick](B)--(2.1,.966667)node[right]{$m\vec v_2$};\draw[dashed](3.8,3.233333)--(3.8,2.1);\draw[dashed](2.1,.966667)--(3.8,2.1);\end{tikzpicture}\end{center}Обоснование.
$1$.  Выберем инерциальную систему отсчёта, связанную с Землёй. За начало отсчёта координат примем первоначальное положение первого шарика.
$2$.  Шарики будем считать материальными точками.
$3$.  Так как сопротивлением воздуха можно пренебречь, то движение шариков можно считать свободным падением.
$4$.  Считаем время взаимодействия шариков при неупругом столкновении малым. Следовательно, импульсом внешней силы (силы тяжести) за это время можно пренебречь. Значит, импульс системы двух шариков при столкновении сохраняется. 

Перейдем к решению.
$1$.  Первый шарик начинает движение из начала координат, а второй  — из точки $A$. До и после столкновения, происходящего в точке $B$, шарики свободно падают. Следовательно, до столкновения для первого шарика

$y_1 ( t ) = v _0y t - \frac{gt ^2}{2} = v _0 \sin \alpha \cdot t - \frac{gt ^2}{2} ,$ 
$v _1y ( t ) = v _0 \sin \alpha - gt,$

а для второго шарика $v _2y ( t ) = - gt.$
$2$.  Шарики сталкиваются в момент $t_{1}$, при этом импульс системы двух шариков сохраняется: $m \vec v _1 + m \vec v _2 = 2m \vec u_0,$ а скорость $\vec u_0$ шариков после удара согласно условию горизонтальна. Следовательно,

\[v _1y ( t_1 ) + v _2y ( t_1 ) = 0 \Longleftrightarrow ( v _0 \sin \alpha - gt_1 ) + ( - gt_1 ) = 0 \Longleftrightarrow t_1 = \frac{v _0 \sin \alpha}{2g} .\]

$3$.  Столкновение шариков происходит на высоте

\[h = y_1 ( t_1 ) = v _0 \sin \alpha \cdot t_1 - \frac{gt_1 ^2}{2} = \frac{v _0 ^2 \sin ^2 \alpha}{2g} - \frac{v _0 ^2 \sin ^2 \alpha}{8g} = \frac{3}{8} \cdot \frac{v _0 ^2 \sin ^2 \alpha}{g} .\]

$4$.  Скорость $\vec u_0$ шариков после удара горизонтальна, поэтому интервал времени $\tau$ от столкновения шариков до их падения на Землю находится из условия $h = \frac{g \tau ^2}{2} ,$ откуда $\tau = \sqrt{\frac{2h}{g}} = \sqrt{3} \cdot \frac{v _0 \sin \alpha}{2g} .$ Таким образом, $v _0 = \frac{2g \tau}{\sqrt{3} \sin \alpha} .$

Ответ: $v _0 = \frac{2g \tau}{\sqrt{3} \sin \alpha} .$
}

\end{document}
